\chapter{Materials and Methods} \label{chap:Materials_and_Methods}

The tools used in this project were installed with the Conda package distribution system version 2-2.4.0 \autocite{anaconda_software_distribution_anaconda_2020}. With Conda, Python packages containing e.~g.~ full software tools can be managed by local environments for simple access without the necessity of a complex system-wide installation. By creation of an environment similar to the one used in this project, with conda environment creation from YAML file, the results can be easily recreated on any modern UNIX-like operating system that has Conda installed. 

\begin{lstlisting}[language=sh]
git clone https://github.com/ahenoch/Masterthesis.git
cd Masterthesis
conda env create -f Environment.yml
\end{lstlisting}  

The YAML file containing the informations for installation like e.~g.~ the used tools versions is present in the \href{https://github.com/ahenoch/Masterthesis.git}{Projects GitHub Repository} an can be made available by e.~g.~ cloning the repository (the important packages and their related versions are also listed in \autoref{tab:Package_Version}).

\begin{table}[!hbt]
    \centering
    \caption[Package Version]{\textbf{Package Version.} The packages that need to be installed by Conda in a specific version for the pipeline to work as expected are listed in this table. Other related packages necessary for execution of the listed ones are meanwhile installed automatically by Conda.}
    \label{tab:Package_Version}
    \pgfplotstabletypeset[
        every head row/.style={
            before row={
                \toprule
            },
            after row={
                \midrule
            },
        },
        every last row/.style={
            after row={
                \bottomrule
            },
        },
        begin table=\begin{tabular*}{.5\textwidth},
        end table=\end{tabular*},
        columns={0,1,2},
        columns/0/.style={string type, multicolumn names=l,column name=\textbf{Name}, column type=@{\extracolsep{\fill}\hspace{6pt}}l},
        columns/1/.style={string type, multicolumn names=l,column name=\textbf{Version}, column type=l},
        columns/2/.style={string type, multicolumn names=l,column name=\textbf{Channel}, column type=l}
    ]
    {Graphics/Packages.csv}
\end{table}

Since its file size exceeds the limits of GitHub, the FASTA file containing all the known high quality sequences of the \gls{IAV}, that is used in this project must be manually retrieved in the latest version \footnote{GenBank Genome Sequence/Annotation Update >= 05/2021} from the \href{https://www.fludb.org/brc/home.spg?decorator=influenza}{Influenza Research Database}.

\begin{table}[!hbt]
    \centering
    \caption[Search Parameter for FASTA file]{\textbf{Search Parameter for FASTA file.} The parameters to use on the nucleotide sequence search interface of the \href{https://www.fludb.org/brc/home.spg?decorator=influenza}{Influenza Research Database}. All paremeters have to be precisely as listed for a exact replication of the FASTA file used in this project.}
    \label{tab:Search}
    \pgfplotstabletypeset[
        every head row/.style={
            before row={
                \toprule
            },
            after row={
                \midrule
            },
        },
        every last row/.style={
            after row={
                %... & ... & ... & ... & ... & ... & ... & ...\\
                \bottomrule
            },
        },
        begin table=\begin{tabular*}{.5\textwidth},
        end table=\end{tabular*},
        columns={0,1},
        columns/0/.style={string type, multicolumn names=l,column name=\textbf{Field}, column type=@{\extracolsep{\fill}\hspace{6pt}}l},
        columns/1/.style={string type, multicolumn names=l,column name=\textbf{Parameter}, column type=l},
    ]
    {Graphics/search.csv}
\end{table}

The FASTA file can be retrieved by navigating from \textit{SEARCH DATA} to \textit{Search Sequences} and then to \textit{Nucleotide Sequences}. On the following page the parameters given in \autoref{tab:Search} have to be used to request the correct file. Following the search request \textit{Select all X segments} have to be checked and after \textit{Download}, \textit{Segment FASTA} has to be checked in \textit{Specify Download Type} setting and \textit{Custom format - select fields from list} in \textit{Format for FASTA file definition line}. Accession Number, Strain Name, Segment, Protein Symbol, Type, SubType, Date, Host Species, Curation Flag have then to be added in the following order. After downloading the FASTA file it has to be placed as \textit{A.fasta} in the local root of the \href{https://github.com/ahenoch/Masterthesis.git}{Projects GitHub Repository}. The version used for the proposed results was acquired at 08/11/2020 \footnote{GenBank Genome Sequence/Annotation Update <= 11/2020 }. Newer versions might change the results slightly.

% \begin{enumerate}[noitemsep]
%     \item \textit{SEARCH DATA}
%     \item \textit{Search Sequences}
%     \item \textit{Nucleotide Sequences}
%     \item \textit{Data Type:} Genome Segments, \textit{Virus Type:} A, \textit{Complete Genome:} Complete Genome Only, \textit{Select Segments:} All, \textit{Complete:} All
%     \item \textit{Search}
%     \item \textit{Select all X segments}
%     \item \textit{Download}
%     \item \textit{Specify Download Type:} Segment FASTA
%     \item \textit{Format for FASTA file definition line:}  Custom format - select fields from list
%     \item \textbf{add in the following order:} Accession Number, Strain Name, Segment, Protein Symbol, Type, SubType, Date, Host Species, Curation Flag
%     \item save as A.fasta in the cloned Masterthesis folder
% \end{enumerate}
%or from the \href{https://github.com/ahenoch/Masterthesis.git}{Projects GitHub Repository}. Using the manually retrieved updated version of the FASTA file might change the recreated results slightly.

To access the Jupyter Notebooks holding the cluster and analysis pipeline, as well as recreate the results, the IPYNB Files in the \href{https://github.com/ahenoch/Masterthesis.git}{Projects GitHub Repository} have to be opened and started in a Jupyter Lab server \autocite{kluyver_jupyter_2016}.

\begin{lstlisting}[language=sh]
conda activate Masterthesis
jupyter lab
\end{lstlisting}  

The raw pipeline clusters \gls{IAV} genomes and saves the results as CSV files and colored \gls{HDBSCAN} clustertrees using ETE3 version 3.1.2 \autocite{huerta-cepas_ete_2016}. It is designed to function as a clustering tool, without the need of a running Jupyter Lab server and with easier modification and faster execution in mind (\autoref{fig:Vectorization_Pipeline} and \autoref{fig:Clustering_Pipeline}). 

\begin{lstlisting}[language=sh]
conda activate Masterthesis
python3 Clustering.py -i A.fasta -o PCA_raw -p 50
\end{lstlisting}  

All settings available to the present day are listed on the next page.

\begin{leftbar}
    \textbf{Clustering.py}
    \begin{nstabbing}
        \qquad\=\qquad\qquad\qquad\qquad\quad\=\kill
    
        -i \> -{}-{}infile \> [path to input file (e.~g.~ A.fasta)]\\
        
        -o \> -{}-{}outfolder \> [path to input file (default: Results)]\\
        
        -s \> -{}-{}segments \> [segments to run the pipeline on (default: 1 2 3 4 5 6 7 8)]\\
        
        -c \> -{}-{}custom\_header \> [if using FASTA with a custom header, every part of it has\\
        
        \> \> to be declared (default: accession strain segment protein\\
        
        \> \> genus subtype date host curation genome)]\\
        
        -m \> -{}-{}metric \> [metric to use in the pipeline (default: cosine)]\\
        
        -mc \> -{}-{}min\_clust \> [min\_cluster\_size parameter for HDBSCAN (default: 2)]\\
        
        -ms \> -{}-{}min\_sample \> [min\_samples parameter for HDBSCAN (default: 1)]\\
        
        -n \> -{}-{}neigbors \> [n\_neighbors parameter for UMAP (default: 100)]\\
        
        -u \> -{}-{}umap \> [n\_components parameter for UMAP (optional)]\\
        
        -p \> -{}-{}pca \> [n\_components parameter for PCA (optional)]
    \end{nstabbing}
\end{leftbar}

\begin{figure}[!hbt]
    \centering
    \includegraphics[width=\textwidth]{Graphics/Vectorization.pdf}
    \caption[Clustering Pipeline]{\textbf{Vectorization Pipeline.} The vectorization pipeline of the proposed clustering tool. The FASTA file is translated to vectors containing the k-mer frequencies of the specific sequence (\autoref{sec:Frequency}). As shown in pathway \textsf{\textbf{1}}, by normalization with L1-norm and \gls{PCA} a low complexity representation of the vectors is obtained for clustering (\autoref{sec:Normalization} and \autoref{sec:PCA}). Pathway \textsf{\textbf{2}} describes additional execution of \gls{UMAP} that can be used after \gls{PCA} as intermediate instead of final step (\autoref{sec:UMAP}).}
    \label{fig:Vectorization_Pipeline}
\end{figure}

\begin{figure}[!hbt]
    \centering
    \includegraphics[width=\textwidth]{Graphics/Clustering.pdf}
    \caption[Clustering Pipeline]{\textbf{Clustering Pipeline.} Following the vectorization pipeline (\autoref{fig:Vectorization_Pipeline}) normalization is used again with L2-norm as preparation for \gls{HDBSCAN} (\autoref{sec:Normalization} and \autoref{sec:HDBSCAN}) Final hybrid clustering (\autoref{sec:HDBSCAN}) is prepared either with the Kneedle Algorithm (pathway \textsf{\textbf{4}} and \autoref{sec:Kneedle}) or DBCV exploration (pathway \textsf{\textbf{3}} and \autoref{sec:HDBSCAN}).}
    \label{fig:Clustering_Pipeline}
\end{figure}

The tool takes around one hour to cluster all eight segments of \gls{IAV} and create the output in form of tables and clustertree graphics with the suggested settings. 

\section{K-Mer Frequency} \label{sec:Frequency}

The FASTA file containing the to be clustered genomes of \gls{IAV} were translated to vectors for clustering in high dimension by counting their 7-mer frequency (\autoref{fig:Vectorization_Pipeline} \textsf{\textbf{A}}). Let $a_n$ be a mathematical sequence of all permutations with repetition of the set $\{A,C,G,T\}$ to the magnitude of seven, thus $4^7$ elements (\autoref{eq:alphabet}) and $s_n$ a mathematical sequence of all genome sequences in the FASTA file (\autoref{eq:alphabet}).

\begin{empheq}{alignat = -1}
    &a_n &&= ( \text{AAAAAAA} , \text{AAAAAAC} , \text{AAAAAAG} , \ldots , \text{TTTTTTT} &&)\label{eq:alphabet}\\
    &s_n &&= ( \text{GCAAAA...} , \text{GCAAAA...} , \text{AGCAAA...} , \ldots , \text{ATGGC...} &&)\label{eq:sequences}
\end{empheq}

The elements of $a_n$ represent all possible constellations of 7-mers. The vector component $x_{i,j}$ of a genome sequence $s_j$ with the length $l$ were calculated by summing up the number of occurrences of element $a_i$ in all the 7-mers of genome sequence $s_j$ (\autoref{eq:frequency}). 

\begin{empheq}{alignat = -1}
    &x_{i,j} = \sum^{l-7+1}_{k=1} \delta(s_i[k\ldots k+7-1], a_j), \ \text{ with }\delta(s,a)=\left\{ \begin{array}{c}1\text{ if \ensuremath{s=a}}\\0\text{ if \ensuremath{s\neq a}}\end{array}\right.\label{eq:frequency}
\end{empheq}

This calculation was repeated for all of the $4^7$ elements of $a_n$ building the vector $\mathbf{x}_i$ (\autoref{eq:full_vector}).

\begin{empheq}{alignat = -1}
    &\mathbf{x}_i = \begin{pmatrix}x_{i,1}\\x_{i,2}\\x_{i,3}\\\vdots\\x_{i,4^7}\end{pmatrix}\label{eq:full_vector}
\end{empheq}

By repeating the calculation \autoref{eq:frequency} for all sequences in $s_n$ a matrix $\mathbf{M}$ was created, holding the k-mer frequency vectors of the genome sequences (\autoref{eq:full_matrix}). The entries in the $j$ columns represent the number of occurencences of 7-mer $a_j$ in a genome sequence $s_i$. The rows represent the $i$ genome sequences from $s_n$. In summary e.g. $\mathbf{M}_{1,1}$ represents the number of occurences of the first 7-mer in $a_n$ in the first genome sequence of $s_n$.

\begin{empheq}{alignat = -1}
    &\mathbf{M} = \begin{bmatrix}\mathbf{x}_1\\\mathbf{x}_2\\\mathbf{x}_3\\\vdots\\\mathbf{x}_i\end{bmatrix}\label{eq:full_matrix}
\end{empheq}

\section{Normalization} \label{sec:Normalization}

The row vectors $\mathbf{x}_i$ of matrix $\mathbf{X}$ were normalized with the L1-norm to align the k-mer frequency vectors located in the rows to an equal length of one (\autoref{fig:Vectorization_Pipeline} \textsf{\textbf{B}}) (\autoref{eq:l1_norm} to \autoref{eq:l1_result}). The \textbf{preprocessing.normalize} function from the \textbf{scikit-learn (sklearn)} package version 0.24.1, with \colorbox{backcolour}{norm='l1'} setting was used for the normalization \autocite{pedregosa_scikit-learn_2011}. Execution of the function on matrix $\mathbf{X}$ results in the matrix $\mathbf{X}_{\text{L1}}$ containing row vectors whose L1-norm $\Vert\mathbf{x}_i\Vert_1$ equals one.

%\begin{lstlisting}[language=python]
%sklearn.preprocessing.normalize(M, norm='l1')
%\end{lstlisting}  

\begin{empheq}{alignat = -1}
    \mathbf{x}_i &= \frac{\mathbf{x}_i}{\Vert\mathbf{x}_i\Vert_1} \label{eq:l1_norm}
\end{empheq}

\begin{empheq}{alignat = -1}
    \Vert\mathbf{x}_i\Vert_1 &= 1\label{eq:l1_result}
\end{empheq}

\begin{empheq}{alignat = -1}
    \mathbf{X}_{\text{L1}} = \text{NORMALIZE}_{\text{L1}}(\mathbf{X})\label{eq:l1_func}
\end{empheq}

The parameters used in this project with settings varying from the default are listed below. All available settings can be fount in the \href{https://scikit-learn.org/stable/modules/generated/sklearn.preprocessing.normalize.html}{API} \autocite{pedregosa_scikit-learn_2011}.

\begin{leftbar}
    \textbf{sklearn.preprocessing.normalize}
    \begin{nstabbing}
        \qquad\qquad\qquad\qquad\qquad\quad\=\kill
        
        X \> [Input matrix to be normalized]\\
        
        norm \> [Norm used for normalization (default: ’l2’)]
        
    \end{nstabbing}
\end{leftbar}

Due to an open \href{https://github.com/scikit-learn-contrib/hdbscan/issues/69}{issue}\footnote{last accessed 02/06/21} in the GitHub Repository of \gls{HDBSCAN} row-wise normalization with \textbf{preprocessing.normalize} was again used with the \colorbox{backcolour}{norm='l2'} setting, right before clustering. (\autoref{eq:l2_norm} to \autoref{eq:l2_result} and \autoref{fig:Clustering_Pipeline} \textsf{\textbf{E}}). It was stated, normalization with L2-norm prior to clustering with euclidean metric setting, chord distance, which is also not directly available, can be used to approximate cosine distance.

\begin{figure}[!hbt]
    \centering
    \begin{subfigure}[b]{0.475\textwidth}
        \caption[Chord]{\textbf{Chord}}
        \label{subfig:Chord}
        \includegraphics[width=\textwidth]{Graphics/Chord.pdf}
    \end{subfigure}
    \hfill
    \begin{subfigure}[b]{0.475\textwidth}
        \caption[L2]{\textbf{L2}}
        \label{subfig:L2}            
        \includegraphics[width=\textwidth]{Graphics/L2.pdf}
    \end{subfigure}
    \caption[Graphical Background of L2 Normalisation]{\textbf{Graphical Background of L2 Normalisation.} The chord can be calculated with the radius similar for both vectors and the information of the angle of two vectors to the origin of the coordinate system. Since this information is not always available like in this proposed project or sometimes hard to calculate for a high number of vectors, the euclidean distance can instead be used with a L2-norm of both vectors equal to $r$ for similar results. To approximate cosine distance, the same calculation of euclidean distance is used with L2-norm equal to $r = 1$. To always obtain similar values of L2-norm equal to 1, L2-norm normalization is used on every vector.}
    \label{fig:L2_Normalisation_Background}
\end{figure}

\begin{empheq}{alignat = -1}
    \mathbf{\hat{x}}_i &= \frac{\mathbf{x}_i}{\Vert\mathbf{x}_i\Vert_2} \label{eq:l2_norm}
\end{empheq}

\begin{empheq}{alignat = -1}
    \Vert\mathbf{\hat{x}}_i\Vert_2 &= 1\label{eq:l2_result}
\end{empheq}

Chord distance $d_{\text{chord}}$ can be calculated when having two vectors, here as an exmaple named as $\mathbf{x}'$ and $\mathbf{y}'$ with the same L2-norm equal to a radius $r$ of a sphere  centered to the origin of the coordinate system and an angle of $\Theta'$ (\autoref{eq:chord} and \autoref{fig:L2_Normalisation_Background}). The euclidean distance $d_{\text{eucl}}$ is equal to the chord distance for the same vectors $\mathbf{x}'$ and $\mathbf{y}'$ with L2-norm of $r$.

\begin{empheq}{alignat = -1}
    &\Vert\mathbf{x}'\Vert_2 = \Vert\mathbf{y}'\Vert_2 = r &&\Rightarrow d_{\text{chord}}(\mathbf{x}',\mathbf{y}') &&= 2 \cdot r \sin \left(\frac{\Theta'}{2}\right) \label{eq:chord}\\
    &&&&&= 2 \cdot \frac{d_{\text{eucl}}(\mathbf{x}',\mathbf{y}')}{2}\\
    &&&&&= d_{\text{eucl}}(\mathbf{x}',\mathbf{y}')\\
    &&&&&= \Vert\mathbf{x}' - \mathbf{y}'\Vert_2
\end{empheq}

Thus in this project chord distance can be calculated with the euclidean distance metric, posterior to the normalization with the L2-norm, which scales the vectors to the unit sphere (\autoref{eq:chord_eucl_1} to \autoref{eq:chord_eucl_2} and \autoref{fig:L2_Normalisation_Background}).

\begin{empheq}{alignat = -1}
    &\Vert\mathbf{\hat{x}}\Vert_2 = \Vert\mathbf{\hat{y}}\Vert_2 = 1 &&\Rightarrow d_{\text{chord}}(\mathbf{\hat{x}},\mathbf{\hat{y}}) &&= d_{\text{eucl}}(\mathbf{\hat{x}},\mathbf{\hat{y}})\label{eq:chord_eucl_1}\\
    &&&&&= \Vert\mathbf{\hat{x}} - \mathbf{\hat{y}}\Vert_2 \label{eq:chord_eucl_2}
\end{empheq}

The used chord distance is proportional with the initially intended to use cosine distance as shown in \autoref{eq:chord_cos_1} to \autoref{eq:chord_cos_7}. Dividing the squared chord distance by 2 results in the cosinus distance of the vectors.

\begin{empheq}{alignat = -1}    
    &\Vert\mathbf{\hat{x}}\Vert_2 = \Vert\mathbf{\hat{y}}\Vert_2 = 1 &&\Rightarrow d_{\text{chord}}(\mathbf{\hat{x}},\mathbf{\hat{y}})^2 &&= \Vert\mathbf{\hat{x}} - \mathbf{\hat{y}}\Vert_2^2\label{eq:chord_cos_1}\\
    &&&&&= (\mathbf{\hat{x}} - \mathbf{\hat{y}})^\top (\mathbf{\hat{x}} - \mathbf{\hat{y}})\label{eq:chord_cos_2}\\
    &&&&&= \mathbf{\hat{x}}^\top \mathbf{\hat{x}} - 2 \mathbf{\hat{x}}^\top \mathbf{\hat{y}} + \mathbf{\hat{y}}^\top \mathbf{\hat{y}}\label{eq:chord_cos_3}\\
    &&&&&= 2 - 2\mathbf{\hat{x}}^\top \mathbf{\hat{y}}\label{eq:chord_cos_4}\\
    &&&&&= 2 - 2 \cos(\Theta)\label{eq:chord_cos_5}\\
    &&&&&= 2 \cdot (1 - \cos(\Theta))\label{eq:chord_cos_6}\\
    &&&&&= 2 \cdot d_{\text{cos}}(\mathbf{x},\mathbf{y})\label{eq:chord_cos_7}
\end{empheq}

 Approximation of cosine distance by normalization with L2-norm followed by euclidean distance calculation is thus a possible and in this project used workaround to overcome the impossibility to use cosine distance.

\section{PCA} \label{sec:PCA}

To handle the complexity of the datasets generated in the project, as well as to simplify them with the least loss of information as possible \gls{PCA} was used (\autoref{fig:Clustering_Pipeline} \textsf{\textbf{C}}) \autocite{pearson_liii_1901}.\gls{PCA} is a statistical technique, used to find a new presentation of the dataset with a lower complexity by variance maximizing, uncorrelated variables, the \glspl{PC}, which are linear functions of the previous ones \autocite{jolliffe_principal_2016}. 

For the \gls{PCA} of matrix $\mathbf{X}$ the \textbf{decomposition.PCA} function from the \textbf{scikit-learn (sklearn)} package, with \colorbox{backcolour}{n\_components=30} setting was used \autocite{pedregosa_scikit-learn_2011}.

The matrix of $\mathbf{X}$ for e.g. segment 4 of \gls{IAV} is of size $56617 \times 4^7$ and the number of components to extract 30 of $4^7$ thereby $\approx 0.18\%$. The size limit of the \textbf{decomposition.PCA} function for calculation with default setting \colorbox{backcolour}{svd\_solver='auto'} is $500 \times 500$ and min. 80\% of the smallest dimension of components to extract. Thus \colorbox{backcolour}{svd\_solver='randomized'} setting of \textbf{decomposition.PCA} is used automatically \autocite{pedregosa_scikit-learn_2011}. Randomized truncated \gls{SVD} is calculated by \textbf{decomposition.PCA} as described in \textcite{halko_finding_2010}.

\autoref{eq:PCA_k} and \autoref{eq:PCA_30} illustrate the use of the \gls{PCA} by the method of \textcite{halko_finding_2010} to reduce the matrix $\mathbf{X}$ with $j=4^7$ dimensions to $k=30$ dimensions \autocite{jolliffe_principal_2016, pedregosa_scikit-learn_2011}. This calculation is a representation of the approach used in \autoref{fig:Clustering_Pipeline} pathway \textsf{\textbf{1}} and the settings for the Jupyter results with the use of \gls{PCA} only \autocite{kluyver_jupyter_2016}.

\begin{empheq}{alignat = -1}
    &\mathbf{X}_{(\text{n\_components})} &&= \text{PCA}(\mathbf{X}, \text{n\_components})\label{eq:PCA_k}\\
    &\mathbf{X}_{(30)} &&= \text{PCA}(\mathbf{X}, 30)\label{eq:PCA_30}
\end{empheq}

The number of components to extract was selected by running \gls{PCA} with different settings for \colorbox{backcolour}{n\_components} with matrix $\mathbf{X}$ of segment 4 and comparing the sum of explained variance. Because of the high increase of explained variance from 10 to 20 and also from 20 to 30 components, a value of 30 was used as default in this project. A fairly small value was used, due to increasing computational effort of \gls{PCA} and also to preserve the usability of Spanning Tree calculation with \gls{HDBSCAN} (\autoref{tab:PCA_Dimension}).

% Let $\mathbf{S}$ be a square covariance matrix of size $i \times i$, calculated on the dataset matrix $\mathbf{X}$ with $i$ $j$-dimensional vectors (size $i \times j)$. 

% \autocite{jolliffe_principal_2016}.

% \begin{empheq}{alignat = -1}
%     x^\ast_{i,j} = x_{i,j} - \bar{x}_j
% \end{empheq}

% \begin{empheq}{alignat = -1}
%     &\mathbf{X}^\ast &&= \mathbf{U} \mathbf{L} \mathbf{A}^\top\\
%     &&&= \mathbf{Y}
% \end{empheq}

% \begin{empheq}{alignat = -1}
%     &\mathbf{X}^{\ast^\top} \mathbf{X}^\ast &&= (\mathbf{U} \mathbf{L} \mathbf{A}^\top)^\top (\mathbf{U} \mathbf{L} \mathbf{A}^\top)\\
%     &&&= \mathbf{A}\mathbf{L}\mathbf{U}^\top\mathbf{U} \mathbf{L} \mathbf{A}^\top\\
%     &&&= \mathbf{A}\mathbf{L}^2\mathbf{A}^\top 
% \end{empheq}

% \begin{empheq}{alignat = -1}
%     \mathbf{Y}_{30} = \mathbf{U}_{30} \mathbf{L}_{30} \mathbf{A}_{30}^\top
% \end{empheq}

The parameters used in this project with settings varying from the default are listed below. All available settings can be fount in the
\href{https://scikit-learn.org/stable/modules/generated/sklearn.decomposition.PCA.html}{API} \autocite{pedregosa_scikit-learn_2011}.

% \begin{empheq}{alignat = -1}
%     &\mathbf{X} = \begin{bmatrix}x_{1,1} & x_{1,2} & x_{1,3} & \dots & x_{1,30}\\
%     x_{2,1} & x_{2,2} & x_{2,3} & \dots & x_{2,30}\\
%     x_{3,1} & x_{3,2} & x_{3,3} & \dots & x_{3,30}\\
%     \vdots & \vdots & \vdots & \ddots & \vdots\\
%     x_{i,1} & x_{i,2} & x_{i,3} & \dots & x_{i,30}
%     \end{bmatrix}\label{eq:full_matrix_2}
% \end{empheq}

\begin{leftbar}
    \textbf{sklearn.decomposition.PCA}
    \begin{nstabbing}
        \qquad\qquad\qquad\qquad\qquad\quad\=\kill

        svd\_solver \> [Strategy used to solve the \gls{SVD} (default: auto)]\\
        n\_components \> [Number of components to extract (default: None)]
    \end{nstabbing}
\end{leftbar}

\section{Uniform Manifold Approximation and Projection} \label{sec:UMAP}

\gls{UMAP} was used for dimension reduction similar to PCA, while aiming to preserve the global structure of the data \autocite{mcinnes_umap_2020}. In it's core, it is similar to \gls{t-SNE} with better run time performance and better structure preservation in the lower dimension and no restrictions for embedding \autocite{maaten_visualizing_2008, mcinnes_umap_2020}. Embedding is, hereby, the presentation of the used dataset in the lower dimension.

The parameters of \gls{UMAP} used match most of the parameters listed in the API under section \glqq \gls{UMAP} enhanced clustering\grqq{}\footnote{\url{https://umap-learn.readthedocs.io/en/latest/clustering.html (last accessed 01/07/2020)}} \autoref{fig:Vectorization_Pipeline} \textsf{\textbf{D}}. The settings are proposed to be used with \gls{UMAP} prior to \gls{HDBSCAN} clustering. Still, a higher dimension to embed the data in \texttt{n\_components=30} was chosen to preserve more information. Instead of visualization in two dimensions, 30 were used because the goal was clustering, not plotting (\autoref{sec:PCA}). The neighbors number was also changed. It is recommended to be set in a range of 1 to 100. Based on the number of sequences used e.~g.~, 56617 for segment 4, the highest recommended setting \texttt{n\_neighbors=100} was used to better preserve the global picture of the data. A small value would make \gls{UMAP} focus on small changes but loose the overall structure which is important when using this magnitude of data points. Also based on the input size \texttt{n\_epochs=200} setting was used automatically.

% \autoref{eq:PCA_100} to \autoref{eq:UMAP_100} denote the use of the \gls{PCA} as shown in \autoref{fig:Vectorization_Pipeline} workflow \textsf{\textbf{2}} and the default setting for Jupyter results with the use of \gls{UMAP} \autocite{pedregosa_scikit-learn_2011, jolliffe_principal_2016, kluyver_jupyter_2016, mcinnes_umap_2020}.

% \begin{empheq}{alignat = -1}
%     &\mathbf{X}_{\text{UMAP}} &&= \text{PCA}(\mathbf{X}_{\text{L1}}, 100)\label{eq:PCA_100}
% \end{empheq}

% The dimension of the dataset was therefore first reduced by the mentioned \gls{PCA} and the simplified result then processed by \gls{UMAP} as described above.

% \begin{empheq}{alignat = -1}
%     &\mathbf{Y}_{\text{UMAP}} &&= \text{UMAP}(\mathbf{X}_{\text{UMAP}}, \text{n\_neighbors}, \text{n\_components}, \text{min\_dist}, \text{n\_epochs})\label{eq:UMAP_d}\\
%     &&&= \text{UMAP}(\mathbf{X}_{\text{UMAP}}, 100, 30, 0.0, 200)\label{eq:UMAP_100}
% \end{empheq}

% Important parameters used in this project, including settings varying from the default, are listed below. All available settings with explanation are listed in the \href{https://umap-learn.readthedocs.io/en/latest/api.html}{API}.

% \begin{leftbar}
%     \textbf{umap.UMAP}
%     \begin{nstabbing}
%         \qquad\qquad\qquad\qquad\qquad\quad\=\kill

%         n\_neighbors \> [number of neighbors (default: 15)]\\
        
%         min\_dist \> [min. package distance of points (default: 0.1)]\\
        
%         n\_components \> [dimension to embed in (default: 2)]\\
        
%         n\_epochs \> [number of epochs for training (default: None)]\\
        
%         metric \> [metric to use (default: 'euclidean')]
%     \end{nstabbing}
% \end{leftbar}

\gls{UMAP} was used posterior to dimension reduction with \gls{PCA}, because of the similarity to \gls{t-SNE}. As explained in the API of \gls{t-SNE}\footnote{\url{https://scikit-learn.org/stable/modules/generated/sklearn.manifold.TSNE.html}}, the dimension should be reduced to a reasonable amount prior to execution to reduce noise. In in section \glqq What is the difference between PCA / UMAP / VAEs?\grqq{} of the \gls{UMAP} API\footnote{\url{https://umap-learn.readthedocs.io/en/latest/faq.html}} a pipeline is proposed, to reduce from high dimension with \gls{PCA}, continue with reduction by \gls{UMAP} and cluster with \gls{HDBSCAN} and is, therefore, used as reference. Reduction with \gls{UMAP} to 30 components posterior to \gls{PCA} with 100 components also provided a comfortable balancing of computational effort of both methods, while preserving $\approx 85\%$ explained variance with \gls{PCA} \autocite{mcinnes_umap_2020}. 

%Since the final dimensionality of the resulting matrix $\mathbf{Y}_{\text{UMAP}}$ either way was $<50$, a reduction with \gls{PCA} to $>50$ could be used while still preserving the usability of Spanning Tree calculation with \gls{HDBSCAN} \autocite{mcinnes_hdbscan_2017}. However, it is indeed also possible to use \gls{UMAP} without prior \gls{PCA}, since there is no limit for input date dimensionality \autocite{mcinnes_umap_2020}. 

\section{HDBSCAN} \label{sec:HDBSCAN}

\gls{HDBSCAN} is an clustering algorithm proposed by \textcite{campello_hierarchical_2015, hutchison_density-based_2013}. It is a novel version of \gls{DBSCAN} but executed with varying values of $\varepsilon$ \autocite{hutchison_density-based_2013}. Thus not one specific threshold is used to define the clusters, but instead clusters of varying densities are extracted based on their stability over epsilon \autocite{mcinnes_hdbscan_2017}. The version of \gls{HDBSCAN} used in this project is the accelerated implementation by the \textbf{hdbscan} package version 0.8.26 of \textcite{mcinnes_accelerated_2017} (\autoref{fig:Clustering_Pipeline} \textsf{\textbf{F}}).

Since there is, to my best knowledge, no submission using \gls{HDBSCAN} for a usecase similar to the one proposed in this project, most settings for \gls{HDBSCAN} were selected by parameter exploration, graphical comparisons and reasonable assumptions. For the selection of all settings the \href{https://hdbscan.readthedocs.io/en/latest/api.html}{API} was used as backup \autocite{mcinnes_hdbscan_2017}. 

To capture, in the best case, even genomes with single \glspl{SNP} at rare appearing positions in their own clusters \colorbox{backcolour}{min\_cluster\_size=2} and \colorbox{backcolour}{min\_samples=1} setting was used as default in the project. These settings are in both cases the smallest values possible, to declare as least points as possible as noise. Setting the \gls{UMAP} parameters to focus on the global structure, while cluster as fine as possible might sound contradictory and is therefore discussed in \autoref{sec:Dimension_Reduction}. The default \colorbox{backcolour}{metric='euclidean'} setting was used for all following executions of \gls{HDBSCAN}, with an exception on the precalculated runs, to approximate the use of cosine distance metric as described in \autoref{sec:Normalization}. Therefore matrix $\mathbf{X}_{\text{PCA}}$ and $\mathbf{Y}_{\text{UMAP}}$ have to be L2-normalized as described in \autoref{sec:Normalization} (\autoref{eq:l2_func_x} and \autoref{eq:l2_func_y}). For calculation of the linkage matrix $\mathbf{L}$ standard \gls{HDBSCAN} without $\varepsilon$ was used.

\autoref{eq:HDB} to \autoref{eq:HDB_link_X} demonstrate the use of \gls{HDBSCAN} to calculate the linkage matrix $\mathbf{L}_{\text{PCA}}$ with matrix $\mathbf{X}_{\text{L2}}$ of \autoref{fig:Clustering_Pipeline} pathway \textsf{\textbf{1}} and was performed in the same way for $\mathbf{L}_{\text{UMAP}}$ with $\mathbf{Y}_{\text{L2}}$ of \autoref{fig:Clustering_Pipeline} pathway \textsf{\textbf{2}} \autocite{mcinnes_hdbscan_2017, gower_minimum_1969}.

\begin{empheq}{alignat = -1}
    &\mathbf{X}_{\text{L2}} &&= \text{NORMALIZE}_{\text{L2}}(\mathbf{X}_{\text{PCA}})\label{eq:l2_func_x}\\
    &\mathbf{Y}_{\text{L2}} &&= \text{NORMALIZE}_{\text{L2}}(\mathbf{Y}_{\text{UMAP}})\label{eq:l2_func_y}
\end{empheq}

\begin{empheq}{alignat = -1}
    &\mathbf{L}_{\text{PCA}} &&= \text{HDBSCAN}_{\text{Link}}(\mathbf{X}_{\text{L2}}, \text{min\_cluster\_size}, \text{min\_samples})\label{eq:HDB}\\
    &&&= \text{HDBSCAN}_{\text{Link}}(\mathbf{X}_{\text{L2}}, 2, 1) \label{eq:HDB_link_X}
\end{empheq}

By combining \gls{HDBSCAN} with \gls{DBSCAN} some of the disadvantages of using either of these methods can be overcome \autocite{mcinnes_hdbscan_2017, moulavi_density-based_2014}. Since \gls{DBSCAN} is a hierarchical clustering tool, it is dependent on a strict threshold for clustering. Points not included in this threshold value are omitted \autocite{ester_density-based_1996, schubert_dbscan_2017}. \gls{HDBSCAN} on the other hand tend to create micro-clusters in areas of high density \autocite{mcinnes_hdbscan_2017}. Using the hybrid approach proposed in \autocite{malzer_hybrid_2020} a threshold value $\varepsilon$ can be used to extract clusters but still use the method of \gls{HDBSCAN} instead of emitting other points. This method is useful when having a small cluster size value while still want to cluster high-density areas together \autocite{malzer_hybrid_2020}.

\begin{figure}[!hbt]
    \centering
    \includegraphics[width=\textwidth]{Graphics/HDB.pdf}
    \caption[]{\textbf{.} low dimension representation of hdbscan manual blabla.}
    \label{fig:HDB}
\end{figure}

\begin{figure}[!hbt]
    \centering
    \includegraphics[width=\textwidth]{Graphics/Hybrid.pdf}
    \caption[]{\textbf{.} low dimensional representation of hybrid hdbscan blabla.}
    \label{fig:Hybrid}
\end{figure}

To find an appropriate value for $\varepsilon$ two different methods were used and compared (\autoref{fig:Clustering_Pipeline} pathway \textsf{\textbf{3}} and \textsf{\textbf{4}}. Using the first method, the \gls{DBCV} exploration in \autoref{fig:Clustering_Pipeline} \textsf{\textbf{G}}, execution of \gls{HDBSCAN} was repeated with different settings for $\varepsilon$ and compared by the \gls{DBCV} to find the value of $\varepsilon$ that maximizes the \gls{DBCV} \autocite{moulavi_density-based_2014}. For these repeated executions of \gls{HDBSCAN}, \colorbox{backcolour}{gen\_min\_span\_tree=True} setting was also necessary, to be able to calculate the \gls{DBCV}, by the minimum spanning tree (\autoref{eq:DBCV_X}) \autocite{moulavi_density-based_2014, gower_minimum_1969}. The exploration was executed for $\mathbf{X}_{\text{L2}}$ and $\mathbf{Y}_{\text{L2}}$ to find the optimal $\varepsilon_{\text{PCA}}$ and $\varepsilon_{\text{UMAP}}$.

\begin{empheq}{alignat = -1}
    &\max_{\substack{0 \leq \varepsilon}} \left(\text{HDBSCAN}_{\text{DBCV}}(\mathbf{X}_{\text{L2}}, 2, 1, \varepsilon)\right) = \text{HDBSCAN}_{\text{DBCV}}(\mathbf{X}_{\text{L2}}, 2, 1, \varepsilon') \Rightarrow \varepsilon' = \varepsilon_{\text{PCA}} \label{eq:DBCV_X}
\end{empheq}

Using the second method the optimal value for $\varepsilon_{\text{UMAP'}}$ and $\varepsilon_{\text{PCA'}}$ were calculated using the Kneedle algorithm \autoref{sec:Kneedle} \autocite{halko_finding_2010}.

With the optimal values for $\varepsilon$ found by \gls{DBCV} exploration and the Kneedle Algorithm, as well as with the matrix for the \gls{UMAP} settings and the one for the \gls{PCA} settings, \gls{HDBSCAN} with the hybrid clustering setting was executed four times. Each execution results in the mathematical sequence of cluster names $N$ related to the sequence of genomic sequences $S$ (\autoref{eq:HDB_cluster_PK} to \autoref{eq:HDB_cluster_UD}) (\autoref{fig:Clustering_Pipeline} \textsf{\textbf{H}}) \autocite{mcinnes_hdbscan_2017, malzer_hybrid_2020}. To sum it up, with the \acrshort{UMAP}/\acrshort{DBCV} method, the cluster name of the first genomic sequence of $S$ is the first element of mathematical sequence $N_{\text{UMAP}}$. 

\begin{empheq}{alignat = -1}
    &N_{\text{PCA}} &&= \text{HDBSCAN}_{\text{Hybrid}}(\mathbf{X}_{\text{L2}}, 2, 1, \varepsilon_{\text{PCA}}) \label{eq:HDB_cluster_PK}\\
    &N_{\text{UMAP}} &&= \text{HDBSCAN}_{\text{Hybrid}}(\mathbf{Y}_{\text{L2}}, 2, 1, \varepsilon_{\text{UMAP}}) \label{eq:HDB_cluster_UK}\\
    &N_{\text{PCA'}} &&= \text{HDBSCAN}_{\text{Hybrid}}(\mathbf{X}_{\text{L2}}, 2, 1, \varepsilon_{\text{PCA'}}) \label{eq:HDB_cluster_PD}\\
    &N_{\text{UMAP'}} &&= \text{HDBSCAN}_{\text{Hybrid}}(\mathbf{Y}_{\text{L2}}, 2, 1, \varepsilon_{\text{UMAP'}}) \label{eq:HDB_cluster_UD}
\end{empheq}

The parameters used in this project with settings varying from the default are listed below. All available settings can be fount in the \href{https://hdbscan.readthedocs.io/en/latest/api.html}{API}.

\begin{leftbar}
    \textbf{hdbscan.HDBSCAN}
    \begin{nstabbing}
        \qquad\qquad\qquad\qquad\qquad\quad\=\kill

        min\_cluster\_size \> [min. size of a cluster (default: 5)]\\
        
        min\_samples \> [conservativeness of the clustering (default: None)]\\
        
        cluster\_selection\_epsilon \> [merge clusters below the threshold (default: 0.0)]\\
        
        gen\_min\_span\_tree \> [generate the minimum spanning tree (default: False)]\\
        
        metric \> [metric to use for clustering (default: euclidean)]
        %alpha \> (default: 1.0)
    \end{nstabbing}
\end{leftbar}

%Clustering
%DBCV mehr als 50dims abkacken bla
%Metrik
%Hybrid selection
%Needle
%Linkage matrix

\section{Kneedle} \label{sec:2.6}

\blindtext

\autocite{satopaa_finding_2011}

\section{Biopython \& Scipy} \label{sec:MAFFT}

The colored cluster trees were created by \textbf{ETE3} with a newick file, generated using a feature request, proposed for the scipy package as described in the issue\footnote{\url{https://github.com/scipy/scipy/issues/8274} (last accessed 06/02/21)} (\autoref{fig:Tree_Pipeline} \textsf{\textbf{J}}).

Clusters centroid vectors were selected by calculating euclidean distance between all the vectors of a cluster to each other. The vector with the minimum mean of overall distance was declared as centroid (\autoref{fig:Tree_Pipeline} \textsf{\textbf{K}}). 

The materials described in the following were only used for the analyses in the \autoref{sec:Clustering_Anomalies} to \autoref{sec:Dimension_Reduction} and are, therefore, not included in the novel \gls{IAV} clustering tool (\autoref{fig:Alignment_Pipeline} and \autoref{fig:Precalc_Pipeline}).

All precalculated trees were also created by using the same mentioned vector distance calculation using cosine distance calculation on the L1-norm normalized k-mer frequency vectors instead (\autoref{fig:Precalc_Pipeline} \textsf{\textbf{F}}). The calculated distances were used for \gls{UPGMA} tree building with tree construction functions of the biopython package. DistanceMatrix function was used with standard settings and DistanceTreeConstructor with \texttt{method='upgma'} (\autoref{fig:Alignment_Pipeline} \textsf{\textbf{O}} and \textsf{\textbf{P}}) \autocite{gower_minimum_1969}. The vectors were also clustered by \gls{HDBSCAN} using the calculated distances with \texttt{metric='precalculated'}. Other settings were used as described in \autoref{sec:HDBSCAN} without hybrid clustering.

\begin{figure}[!hbt]
    \centering
    \includegraphics[width=\textwidth]{Graphics/Precalculated.pdf}
    \caption[Precalculation Pipeline]{\textbf{Precalculation Pipeline.} For the precalculated trees the L1-normalized k-mer vectors were compared by cosine distance prior to clustering by \gls{HDBSCAN} and on the other hand processing with \textbf{biopython (Bio)} package \gls{UPGMA} calculation and visualization by \textbf{ETE3}.}
    \label{fig:Precalc_Pipeline}
\end{figure}

MAFFT was used for \glspl{MSA} and guide tree creation with \texttt{treeout=True} on the centroid sequences and on a small FASTA subset (\autoref{fig:Alignment_Pipeline} \textsf{\textbf{L}}) \autocite{katoh_mafft_2013}. The \gls{MSA} of the FASTA subset was converted to evolutionary distances with DistanceCalculator of the biopython package using standard settings and clustered by \gls{HDBSCAN} similar to the cosine distances (\autoref{fig:Alignment_Pipeline} \textsf{\textbf{M}} and \textsf{\textbf{N}}). PairwiseAligner also from the biopython packages was used for the pairwise alignments in \autoref{sec:Serotype_Classification}.

\begin{figure}[!hbt]
    \centering
    \includegraphics[width=\textwidth]{Graphics/Alignment.pdf}
    \caption[Alignment Pipeline]{\textbf{Alignment Pipeline.} The analysis steps involving \glspl{MSA} are performed as pictured. The sequences related to the centroids were aligned using MAFFT resulting in the output as guide-tree visualized by \textbf{ETE3}. The small subset for comparison was also aligned by MAFFT prior to evolutionary distance calculation and clustering by \gls{HDBSCAN}.}
    \label{fig:Alignment_Pipeline}
\end{figure}






For the pairwise alignments in \autoref{sec:Serotype_Classification} from the \textbf{biopython (Bio)} package version 1.78, the \textbf{Align.PairwiseAligner} function was used with standard settings \autocite{cock_biopython_2009}. 

Important parameters used in this project, including settings varying from the default, are listed below. All available settings with explanation are listed in the \href{https://biopython.org/docs/1.75/api/Bio.Align.html}{API}.

The Biopython implementation of MAFFT version 7.475 from the \textbf{biopython (Bio)} package version 1.78 was used to build the guide trees in \autoref{sec:Clustering_Anomalies} according to \autoref{fig:Alignment_Pipeline} pathway \textsf{\textbf{5}} with the genomic sequences of the centroids $C$ \autocite{katoh_mafft_2013, cock_biopython_2009} (\autoref{fig:Alignment_Pipeline}. The settings were the default settings proposed in the \href{https://mafft.cbrc.jp/alignment/software/}{manual}, with changes in the thread number and the output format. For faster execution \colorbox{backcolour}{thread=6} and for export of the guide tree \colorbox{backcolour}{treeout=True} setting was used \autocite{katoh_mafft_2013, cock_biopython_2009}. The guide-tree was colored with \textbf{ETE3} (\autoref{fig:Alignment_Pipeline} \textsf{\textbf{K}}) \autocite{huerta-cepas_ete_2016}.

Important parameters used in this project, including settings varying from the default, are listed below. All available settings with explanation are listed in the \href{https://mafft.cbrc.jp/alignment/software/}{manual}.

% \begin{leftbar}
%     %\textbf{mafft}
%     \textbf{Bio.Align.Applications.MafftCommandline}
%     \begin{nstabbing}
%         \qquad\qquad\qquad\qquad\qquad\quad\=\kill
    
%         treeout \> [export guide tree used for alignment (default: off)]\\
        
%         thread \> [number of used threads (default: 1)]\\
        
%         input \> [input FASTA file]
        
%         %-{}-{}6merpair \> (default: on)\\
%         %\> input\\
%         %> \> output
%     \end{nstabbing}
% \end{leftbar}

In \autoref{sec:Comparison_Clustering} five different cluster trees are compared to each other. The trees are based on clustering with \gls{HDBSCAN}, without hybrid clustering ($\varepsilon=0$) \autocite{malzer_hybrid_2020, mcinnes_hdbscan_2017}. For each clustering a different matrix was used as input. The trees of \gls{UMAP} and \gls{PCA} (\autoref{fig:Simple_Clustertree_PCA} and \autoref{fig:Simple_Clustertree_UMAP}) were created according to \autoref{fig:Vectorization_Pipeline} and \autoref{fig:Clustering_Pipeline} and as described in \autoref{sec:Frequency} to \autoref{sec:HDBSCAN} with a smaller FASTA subset $S$ containing only segment 4 sequences with subtypes H13 and H16 and a resulting smaller L2-norm normalized \gls{PCA} matrix $\mathbf{X}_{\text{L2}}$ and smaller \gls{UMAP} matrix $\mathbf{Y}_{\text{L2}}$. Clustering for the trees of \gls{UMAP} and \gls{PCA} was performed similar to \autoref{sec:HDBSCAN} but without $\varepsilon$ with the small matrix variants of $\mathbf{X}_{\text{L2}}$ and $\mathbf{Y}_{\text{L2}}$ (\autoref{eq:hdb_prime_x} and \autoref{eq:hdb_prime_y}). 

% \begin{empheq}{alignat = -1}
%     &N_{\text{UMAP}} &&= \text{HDBSCAN} (\mathbf{X}_{\text{L2}}, 2, 1, 0)\label{eq:hdb_prime_x}\\
%     &N_{\text{PCA}} &&= \text{HDBSCAN} (\mathbf{Y}_{\text{L2}}, 2, 1, 0)\label{eq:hdb_prime_y}
% \end{empheq}

The two trees based the clustering with precalculation (\autoref{fig:Simple_Clustertree_Cosine} and \autoref{fig:Simple_Clustertree_Euclid}) were created by distance calculation on the L1-norm normalized matrix $\mathbf{X}_{\text{L1}}$ of the smaller $S$. For the calculation of the distance the \textbf{scipy} package version 1.6.0 was used with the \textbf{spatial.distance.cdist} function (\autoref{fig:Precalc_Pipeline} \textsf{\textbf{N}}) \autocite{scipy_10_contributors_scipy_2020}. For calculation of the cosine distance \colorbox{backcolour}{metric='cosine'} and for euclidean distance \colorbox{backcolour}{metric='euclidean'} setting was used (\autoref{eq:c_calc_matrix} to \autoref{eq:e_calc_matrix}). The resulting distance matrices $\mathbf{C}$ and $\mathbf{E}$ were clustered with \gls{HDBSCAN} and \colorbox{backcolour}{metric='precalculated'} setting (\autoref{fig:Precalc_Pipeline} \textsf{\textbf{F}} and \autoref{eq:hdb_prime_e} to \autoref{eq:hdb_prime_c}) \autocite{mcinnes_hdbscan_2017}.

% \begin{empheq}{alignat = -1}
%     &c_{i,j} &&= d_{\text{cos}}(\mathbf{x}_i, \mathbf{x}_j), \ i = 1, \ldots o, \ j = 1, \ldots o\label{eq:c_calc_1}\\
%     &&&= 1 - \frac{\mathbf{x}_i^\top\mathbf{x}_j}{\Vert\mathbf{x}_i\Vert \cdot \Vert\mathbf{x}_j\Vert}\label{eq:c_calc_2}
% \end{empheq}

% \begin{empheq}{alignat = -1}
%     &\mathbf{C} &&= \text{CDIST}_{\text{cosine}}(\mathbf{X}_{\text{L1}}, \mathbf{X}_{\text{L1}}) \label{eq:c_calc_matrix}\\
%     &\mathbf{E} &&= \text{CDIST}_{\text{euclidean}}(\mathbf{X}_{\text{L1}}, \mathbf{X}_{\text{L1}}) \label{eq:e_calc_matrix}
% \end{empheq}

% \begin{empheq}{alignat = -1}
%     &N_{\text{eucl}} &&= \text{HDBSCAN} (\mathbf{E}, 2, 1, 0)\label{eq:hdb_prime_e}\\
%     &N_{\text{cos}} &&= \text{HDBSCAN} (\mathbf{C}, 2, 1, 0)\label{eq:hdb_prime_c}
% \end{empheq}

% Important parameters used in this project, including settings varying from the default, are listed below. All available settings with explanation are listed in the \href{https://docs.scipy.org/doc/scipy/reference/generated/scipy.spatial.distance.cdist.html}{API}.

% \begin{leftbar}
%     \textbf{scipy.spatial.distance.cdist}
%     \begin{nstabbing}
%         \qquad\qquad\qquad\qquad\qquad\quad\=\kill
    
%         metric \> [Distance metric for calculation (default: 'euclidean')]

%     \end{nstabbing}
% \end{leftbar}

The \textbf{spatial.distance.cdist} function is also used for a fast pairwise alignment using the mentioned \textbf{Align.PairwiseAligner} above. Therefore, a small python definition containing the \textbf{Align.PairwiseAligner} function was inserted as metric setting into \textbf{spatial.distance.cdist}. 

The last of the five trees is based on MAFFT with the \textbf{biopython} package. The alignment was done on the subset $S$. On the alignment the distance matrix was calculated by the \textbf{Phylo.TreeConstruction.DistanceCalculator} function from the \textbf{biopython (Bio)} package. The standard settings for nucleotide sequences was used for the calculation \colorbox{backcolour}{model='identity'} \autocite{cock_biopython_2009}.

% \begin{empheq}{alignat = -1}
%     &\mathbf{G} &&= \text{DistanceCalculator}_{\text{identity}}(\text{MAFFT}(S_{\text{Sub}}))
% \end{empheq}

% \begin{empheq}{alignat = -1}
%     &N_{\text{msa}} &&= \text{HDBSCAN} (\mathbf{G}, 2, 1, 0)\label{eq:hdb_prime_g}
% \end{empheq}

% Important parameters used in this project, including settings varying from the default, are listed below. All available settings with explanation are listed in the \href{https://biopython.org/docs/latest/api/Bio.Phylo.TreeConstruction.html}{API}.

% \begin{leftbar}
%     \textbf{Bio.Phylo.TreeConstruction.DistanceCalculator}
%     \begin{nstabbing}
%         \qquad\qquad\qquad\qquad\qquad\quad\=\kill
    
%         model \> [model for distance calculation (default: 'identity')]

%     \end{nstabbing}
% \end{leftbar}

The \textbf{Phylo.TreeConstruction.DistanceMatrix} and \textbf{DistanceTreeConstructor} function, also from the \textbf{biopython} package, was used for calculation of precalculated UPGMA trees in \autoref{sec:K_mer_Representation}. Instead of the default neighbor joining setting, the bottom-up hierarchical clustering method \colorbox{backcolour}{method='upgma'} was used (\autoref{fig:Alignment_Pipeline} \textsf{\textbf{O}} and \textsf{\textbf{P}}) \autocite{gower_minimum_1969, cock_biopython_2009}. The matrices calculated in \autoref{eq:c_calc_matrix} and \autoref{eq:e_calc_matrix} were used as input for the calculation of the UPGMA-trees \autocite{sokal_statistical_1958}.

% Important parameters used in this project, including settings varying from the default, are listed below. All available settings with explanation are listed in the same \href{https://biopython.org/docs/latest/api/Bio.Phylo.TreeConstruction.html}{API} of \textbf{Bio.Phylo.TreeConstruction}.

% \begin{leftbar}
%     \textbf{Bio.Phylo.TreeConstruction.DistanceMatrix}
%     \begin{nstabbing}
%         \qquad\qquad\qquad\qquad\qquad\quad\=\kill
    
%         names \> [name of columns and rows]\\
        
%         matrix \> [input matrix]
%     \end{nstabbing}
% \end{leftbar}

% \begin{leftbar}
%     \textbf{Bio.Phylo.TreeConstruction.DistanceTreeConstructor}
%     \begin{nstabbing}
%         \qquad\qquad\qquad\qquad\qquad\quad\=\kill
    
%         method \> [construction method for the distance tree (default: 'nj')]
        
%     \end{nstabbing}
% \end{leftbar}